% ── 1 Introdução ────────────────────────────────────────────────────────────
\chapter{INTRODUÇÃO}

A administração de uma escola pública envolve tarefas que vão muito além do pedagógico. Coordenações e secretarias lidam diariamente com matrículas, distribuição de turmas, controle de frequência docente e remanejamento de horários. Entre essas demandas, a gestão de ausências de professores se destaca pela imprevisibilidade: atestados médicos, licenças e afastamentos ocorrem sem aviso prévio e exigem resposta imediata para que turmas não fiquem ociosas. Lück (2009) aponta que a eficácia da gestão escolar depende da capacidade de articular processos administrativos e pedagógicos de forma integrada, o que se torna inviável quando ferramentas adequadas de organização não estão disponíveis.

Apesar do avanço da tecnologia em diversas áreas da educação, muitas escolas públicas municipais ainda dependem de planilhas avulsas, cadernos de anotação e comunicação por aplicativos de mensagem para resolver questões operacionais. Moran (2015) argumenta que a incorporação de tecnologias digitais na educação precisa ultrapassar o uso em sala de aula e alcançar também os processos de gestão, possibilitando decisões mais ágeis e fundamentadas em dados. Na prática, a ausência de sistemas informatizados para tarefas administrativas resulta em retrabalho, erros de comunicação e perda de informações relevantes.

Esse cenário se reproduz na Escola Municipal Patrícia Viviani, localizada no bairro Topolândia, em São Sebastião (SP). A escola funciona com múltiplas turmas e turnos, e o controle de substituições de professores é feito manualmente pela secretaria e pela vice-direção. Quando um professor comunica ausência, a equipe administrativa precisa verificar quais horários ficaram descobertos, identificar substitutos disponíveis e reorganizar a grade, tudo isso frequentemente no mesmo dia. Choques de horários, substituições duplicadas e a falta de registro centralizado são problemas recorrentes. A Vice-Diretora Marília Izabel confirmou que a escola não dispõe de nenhum sistema digital voltado a esse tipo de controle, o que torna o processo lento e sujeito a falhas.

No semestre anterior, durante o Projeto Integrador II, o grupo desenvolveu um protótipo do sistema Leciona para atender a essa demanda. Construído com o framework Django e banco de dados PostgreSQL, o protótipo permitia o cadastro de professores, disciplinas, turmas e horários, além de registrar ausências e sugerir substituições. No entanto, o sistema funcionava apenas em ambiente local, sem possibilidade de acesso remoto por parte da equipe escolar. Não havia testes automatizados que garantissem a integridade das regras de negócio (como a impossibilidade de alocar um mesmo professor em duas turmas no mesmo horário), e a interface, embora funcional, não oferecia painéis analíticos para apoiar a tomada de decisão da coordenação.

A continuidade do projeto no Projeto Integrador III se justifica tanto pela demanda real da escola parceira quanto pelo alinhamento com o tema norteador do semestre, que exige desenvolvimento de software web em nuvem com banco de dados, automação por script web (JavaScript), acessibilidade, integração contínua, testes automatizados e fornecimento de API RESTful. Manter o projeto anterior, em vez de iniciar um novo, permite ao grupo aprofundar a solução: corrigir limitações identificadas na primeira versão, incorporar boas práticas de engenharia de software e entregar um produto com maior maturidade técnica. Sommerville (2011) destaca que a evolução de software é uma etapa natural e necessária do ciclo de vida de sistemas, na qual requisitos são refinados e a arquitetura é ajustada conforme o conhecimento adquirido durante o uso e o desenvolvimento.

A decisão de evoluir o Leciona também considerou o retorno da escola. Após a apresentação do protótipo no PI2, a Vice-Diretora indicou funcionalidades prioritárias, como a visualização consolidada de faltas por período e a comunicação antecipada de substituições aos professores. Essas demandas reforçaram a necessidade de uma API que permitisse o consumo dos dados por diferentes interfaces, além do painel web já existente.

Este relatório parcial apresenta a evolução do sistema Leciona no contexto do Projeto Integrador III. O foco recai sobre os módulos de gestão de horários, registro de ausências, alocação de substituições e comunicação com a equipe escolar. São descritas a arquitetura modular adotada, a estratégia de deploy contínuo em nuvem utilizando Docker, a implementação de testes automatizados para validação de regras de negócio, a construção da API RESTful e o desenvolvimento de um dashboard analítico para suporte à gestão. O documento também detalha a metodologia empregada, o cronograma de atividades e os resultados parciais obtidos até o momento.
